\documentclass[11pt]{article}
\usepackage[T1]{fontenc}
\usepackage[utf8]{inputenc}
\usepackage{amsmath,amssymb,amsthm,mathtools}
\usepackage[margin=1in]{geometry}
\usepackage{booktabs}
\usepackage{array}
\usepackage{microtype}
\usepackage{enumitem}
\usepackage[font=small,labelfont=bf]{caption}
\usepackage{xcolor}
\usepackage{tikz}
\usepackage{listings}
\setlength{\emergencystretch}{2.5em}
\usepackage[obeyspaces,spaces]{url} % keep the spaces inside \lean{decide +kernel}, \lean{lake build}, ...
\usepackage{hyperref}
\hypersetup{colorlinks=true, linkcolor=black, citecolor=black, urlcolor=blue}
\usepackage[capitalise,noabbrev]{cleveref}
% Allow URLs (and Lean names) to break at hyphens as well.
\makeatletter\g@addto@macro\UrlBreaks{\do\-}\makeatother

\newtheorem{theorem}{Theorem}[section]
\newtheorem{proposition}[theorem]{Proposition}
\newtheorem{lemma}[theorem]{Lemma}
\newtheorem{corollary}[theorem]{Corollary}
\theoremstyle{definition}
\newtheorem{definition}[theorem]{Definition}
\newtheorem{remark}[theorem]{Remark}
% cleveref names shared-counter environments after the counter ("Theorem"); alias them per environment.
\usepackage{etoolbox}
\AtBeginEnvironment{lemma}{\crefalias{theorem}{lemma}}
\AtBeginEnvironment{proposition}{\crefalias{theorem}{proposition}}
\AtBeginEnvironment{corollary}{\crefalias{theorem}{corollary}}
\AtBeginEnvironment{definition}{\crefalias{theorem}{definition}}
\AtBeginEnvironment{remark}{\crefalias{theorem}{remark}}

% Lean identifiers: typewriter, with line breaks allowed at dots and underscores.
\DeclareUrlCommand\lean{\urlstyle{tt}}
\newcommand{\re}{\operatorname{Re}}
\newcommand{\im}{\operatorname{Im}}
\newcommand{\C}{\mathbb{C}}
\newcommand{\R}{\mathbb{R}}
\newcommand{\N}{\mathbb{N}}
\newcommand{\Z}{\mathbb{Z}}
\newcommand{\Q}{\mathbb{Q}}
\newcommand{\OK}{\mathcal{O}_K}
\newcommand{\disc}{\operatorname{disc}}
\newcommand{\Nm}{\operatorname{N}}
\newcommand{\Cl}{\operatorname{Cl}}

\definecolor{leanbg}{HTML}{F4F2EC}

% The formal side of each step, set off from the mathematics.
\newenvironment{inlean}{\par\medskip\noindent\begin{tikzpicture}\node[fill=leanbg,rounded corners=2pt,
  inner sep=6pt,text width=\dimexpr\linewidth-12pt\relax,font=\small]\bgroup\textbf{In Lean.}\ }%
  {\egroup;\end{tikzpicture}\par\medskip}

\lstdefinelanguage{Lean}{
  morekeywords={theorem,lemma,def,noncomputable,import,namespace,end,open,by,fun,sorry,example,variable,structure,where},
  sensitive=true, morecomment=[l]{--}, morecomment=[s]{/-}{-/}, morestring=[b]",
}
\lstset{language=Lean, basicstyle=\small\ttfamily, backgroundcolor=\color{leanbg},
  frame=single, framerule=0pt, framesep=5pt, xleftmargin=5pt, xrightmargin=5pt,
  breaklines=true, columns=fullflexible, keepspaces=true, showstringspaces=false,
  literate={ℝ}{{$\mathbb{R}$}}1 {ℂ}{{$\mathbb{C}$}}1 {ℕ}{{$\mathbb{N}$}}1 {ℤ}{{$\mathbb{Z}$}}1 {ℚ}{{$\mathbb{Q}$}}1
           {∀}{{$\forall$}}1 {∃}{{$\exists$}}1 {→}{{$\to$}}1 {≠}{{$\neq$}}1 {≤}{{$\le$}}1 {≥}{{$\ge$}}1
           {∧}{{$\wedge$}}1 {∨}{{$\vee$}}1 {¬}{{$\neg$}}1 {↔}{{$\leftrightarrow$}}1 {∈}{{$\in$}}1 {∉}{{$\notin$}}1
           {α}{{$\alpha$}}1 {χ}{{$\chi$}}1 {σ}{{$\sigma$}}1 {ω}{{$\omega$}}1 {θ}{{$\theta$}}1 {π}{{$\pi$}}1
           {∑}{{$\sum$}}1 {↦}{{$\mapsto$}}1 {·}{{$\cdot$}}1 {⟨}{{$\langle$}}1 {⟩}{{$\rangle$}}1
           {₀}{{$_0$}}1 {²}{{$^2$}}1 {ε}{{$\varepsilon$}}1 {𝓞}{{$\mathcal{O}$}}1 {∣}{{$\mid$}}1
           {⊔}{{$\sqcup$}}1 {Γ}{{$\Gamma$}}1 {√}{{$\surd$}}1 {′}{{$'$}}1 {ζ}{{$\zeta$}}1
           {⁴}{{$^4$}}1 {⁵}{{$^5$}}1 {⌊}{{$\lfloor$}}1 {⌋}{{$\rfloor$}}1 {₊}{{$_+$}}1
           {⁻}{{$^{-}$}}1 {¹}{{$^{1}$}}1 {⁰}{{$^{0}$}}1 {‖}{{$\Vert$}}1 {⊤}{{$\top$}}1 {⊢}{{$\vdash$}}1}


\title{Tenth-discriminant bounds compared:\\ a verified survey of the methods in Tao's challenge,\\ and an analytic cutoff $q < 565$ from the determinant method}
\author{Peter W. Murphy}
\date{October 2026 (v1.3.0)\\[2pt] \small Repository \url{https://github.com/DrMurphyIsIn/tenth-discriminant}, directory \lean{comparison/}\\ \small DOI \url{https://doi.org/10.5281/zenodo.23665245}\\ \small (all versions: \url{https://doi.org/10.5281/zenodo.23550257})\\ \small \textbf{Status:} certified by interval arithmetic; reproduced by a separately written implementation; \textbf{not kernel-checked}}

\begin{document}
\maketitle

\begin{abstract}
Terence Tao's tenth-discriminant challenge asks, for each recent argument excluding Landau--Siegel zeros or establishing a quasi-Riemann hypothesis, for the effective upper bound $D_0$ it gives on a hypothetical tenth imaginary quadratic field of class number one, so that the strength of the arguments can be compared. The comment thread now holds about twenty such extractions. We give, first, a survey that puts them side by side: each value is quoted at source with its author and comment, and labelled by what we could check. For every input that can be written as a real-zero-free segment we use one reduction, the Chowla--Selberg expansion of the principal form, and certify in ball arithmetic: $d \le 652$ from the $7/8$ half-plane, $d \le 1151$ from the $11/12$ half-plane (LemmaLab's value, comment 698473), $d \le 199$ under GRH, with the limit $64\pi^2 e^{-2\gamma} = 199.1208\ldots$, $d < 10^{10}$ from Banks's explicit theorem (as twinheterodox and michaelmross posted), and, from Kimura's zero-gap theorems (Lean-formalised by the author; not replayed by us), $d < 6.47\cdot 10^7$ (his chain theorem 58) and $d < 1.23\cdot 10^8$ (the input of Rudolph's extraction). Reading the same curve backwards puts every method on one scale, the zero-free abscissa it is ``worth'': the light-machinery determinant bound $q<570$ of twinheterodox corresponds to $0.861$, ahead of the $7/8$ half-plane on this problem. Second, we rebuild that determinant bound (v0.6.1 of twinheterodox, after Green, van Dijk, Pratt and Kerr) from our own derivation: two separately written AI implementations reproduce its certificate row by row, to $10^{-13}$, in its own configuration, and the method passes a truth test on the fields of discriminant $-163$, $-67$ and $-43$. Third, we improve it: an explicit treatment of the Poisson modes in the norm series, the sharper Rosser--Schoenfeld constant, re-optimised sieve rows and kernels, and a band-uniformity check by Taylor models. Together with the kernel-checked bound $d \le 652$ of our earlier work, this gives the analytic cutoff $h(K) = 1 \Rightarrow q < 565$. The improvement is about $0.9\%$ and changes no candidate discriminant; it is a measure of the method's strength, which is what the challenge compares. The computations are certified by interval arithmetic and reproduced by an independent implementation; they are \emph{not} kernel-checked.
\end{abstract}

%=====================================================================================================
\section{Introduction}
%=====================================================================================================

\subsection{The challenge}
There are nine imaginary quadratic fields of class number one, with discriminants $-3$, $-4$, $-7$, $-8$, $-11$, $-19$, $-43$, $-67$, $-163$. That there is no tenth was proved by Heegner \cite{Heegner1952}, with later proofs by Baker \cite{Baker1966} and Stark \cite{Stark1967}. Tao's challenge \cite{Tao2026} takes the recent effective arguments against Landau--Siegel zeros, and the recent quasi-Riemann hypotheses, and asks what bound each gives on a \emph{hypothetical} tenth discriminant: ``if a tenth existed, then $d \le D_0$''. The answer is known; the numbers measure the arguments. Tao asks for every approach's value, ``even if they are worse'', and for negative results.

\subsection{What this paper does}
\begin{enumerate}[leftmargin=*]
\item \textbf{A survey} (\cref{sec:survey}). Every posted extraction, quoted at source with comment number and author, in one table, with a status per row: reproduced independently by us, cited, or not reproducible. For the zero-free inputs we use one reduction (\cref{sec:reduction}) so that rows compare methods rather than reductions, and we read the resulting curve backwards to obtain a single strength scale (\cref{sec:scale}).
\item \textbf{A reproduction} (\cref{sec:reproduction}) of the strongest posted light-machinery bound, $q < 570$ \cite{twinheterodox}, from our own derivation and code.
\item \textbf{An improvement} (\cref{sec:improvement,sec:theorem}) of that bound to $q < 565$.
\end{enumerate}

\subsection{How the claim is phrased}\label{sec:framing}
The theorem of \cref{sec:theorem} is an \emph{analytic cutoff}: an argument that excludes every $q \ge 565$ without a finite search above $565$. The finite range below the cutoff is checked by computation (\cref{sec:finite}). Above $652$ the argument is our kernel-checked theorem $d \le 652$ \cite{CN1paper}, which uses OpenAI's $7/8$ half-plane \cite{OAIqrh}; between $565$ and $652$ it is the determinant method alone. So the result belongs to the class ``unconditional, Heegner-free, with a quasi-Riemann input''. Since a light-machinery bound is also valid in that class, the comparison with $570$ is made inside it.

\textbf{What is verified, and how.} Every inequality of the certificate is evaluated in Arb ball arithmetic \cite{Arb}. A second implementation, written separately by an independent AI review session in high-precision floating point, reproduces the certificate band by band. Both implementations are AI-written; no human has checked them. \textbf{Nothing in \cref{sec:reproduction,sec:improvement,sec:theorem} is kernel-checked.} The reason is given in \cref{sec:limits}.

%=====================================================================================================
\section{Conventions}\label{sec:conventions}
%=====================================================================================================
Throughout, $D<0$ is a fundamental discriminant, $q = d = |D|$, $K = \Q(\sqrt D)$, $\chi = \chi_D$. All bounds $D_0$ are upper bounds: $h(K) = 1 \Rightarrow q \le D_0$ (or $< D_0$, as stated). Lower bounds on a hypothetical tenth discriminant, such as $q > \exp(2.2\cdot10^7)$, which Pratt (comment 698460) attributes to Stark's 1966 computation of low-lying zeros, are finite-range tools and are not compared here.

Two elementary consequences of $h(K) = 1$ are used repeatedly; they are classical, and were emphasised in the thread by Silberman (comment 698449) and Pratt (comment 698460). Every prime $p < q/4$ is inert, because a non-rational element of $\OK$ has norm at least $q/4$. Hence, for $q > 8$, $q$ is a prime and $q \equiv 3 \pmod 8$.

%=====================================================================================================
\section{One reduction for the zero-free inputs}\label{sec:reduction}
%=====================================================================================================
For $h(K) = 1$ and $q > 4$, $\zeta_K(s) = \tfrac12 Z_Q(s)$ for the principal form $Q$, and Chowla--Selberg \cite{ChowlaSelberg1967} gives
\begin{equation}\label{eq:cs}
\zeta_K(s) = \zeta(2s) + 2^{2s-1}\sqrt\pi\,\frac{\Gamma(s-\frac12)}{\Gamma(s)}\,\zeta(2s-1)\,q^{\frac12 - s} + E(s),
\end{equation}
with the Bessel term $E$ bounded by $2^{s+3/2}\pi^s\Gamma(s)^{-1}q^{1/4-s/2}\sum_n n^{s-1/2}\sigma_{1-2s}(n)K_{s-1/2}(\pi n\sqrt q)$. We checked the normalisation against direct lattice sums before using it. In the thread this expansion was written down by michaelmross (comment 698466) and evaluated at $s = 7/8$ by LemmaLab (comment 698472). Write $M(\sigma, q)$ for the sum of the first two terms and $d_M(\sigma)$ for the root in $q$ of $M(\sigma, q) = 0$, which has the closed form $d_M(\sigma) = \bigl(-C(\sigma)\zeta(2\sigma-1)/\zeta(2\sigma)\bigr)^{1/(\sigma-1/2)}$.

\begin{proposition}[the edge form]\label{prop:edge}
If $L(s,\chi)$ has no real zero in $(\sigma_0, 1)$, $\tfrac12 < \sigma_0 < 1$, and $h(K) = 1$, then $q < d_M(\sigma_0)$ (at $\sigma_0 = 7/8$ this gives $d \le 652$, since $-653$ is not a discriminant: $653 \equiv 1 \pmod 4$). The certificate is $M(\sigma_0, D) - B(\sigma_0, D) > 0$ at $D = \lceil d_M(\sigma_0)\rceil$, with $M$ increasing and $B$ decreasing in $q$.
\end{proposition}

\begin{proposition}[zero forcing; twinheterodox, comment 698456, and michaelmross, comment 698466]\label{prop:zero}
If $h(K) = 1$ and $q \ge 10^{10}$, then $L(s,\chi)$ has a real zero $\beta$ with $1 - \beta < 1.92/\sqrt q$. More generally $\kappa(q_1)$ is certified for $q \ge q_1$, with $\kappa \to 6/\pi$.
\end{proposition}
The route is twinheterodox's (a theta integral, with $1-\beta \le 10/\sqrt q$) and, independently, michaelmross's (the Chowla--Selberg expansion, with $2/\sqrt d$ and the asymptotic constant $6/\pi$); the constant $1.92$ for $q \ge 10^{10}$ is ours. The proof uses $\zeta(u) > 1/(u-1) + \frac12$ on $(0,1)$ (first-order Euler--Maclaurin) and the sign change of $\zeta_K$ near $s = 1$. For an input in the normal form ``$(1-\beta)\log q \ge c$ for $q \ge q_0$'', \cref{prop:zero} gives $D_0 = \max(q_0, d_c)$, with $d_c$ the point where $\kappa\log x/\sqrt x$ falls below $c$; for a piecewise constant $c(\log q)$, $D_0$ is the least $d$ beyond which $\kappa\log x/\sqrt x < c(\log x)$ for all $x$.

%=====================================================================================================
\section{The survey}\label{sec:survey}
%=====================================================================================================

\subsection{The table}
\Cref{tab:survey} lists the extractions in the thread of \cite{Tao2026}, read from the page on 25 October 2026 (74 comments), with the formulas taken from the page's image text. Status: \textbf{R} = reproduced independently by us (and by the review session's own code); \textbf{C} = cited as posted, not yet checked by us; \textbf{NR} = not reproducible from what is posted.

\begin{table}[ht]\centering\small
\begin{tabular}{@{}p{3.7cm}p{3.6cm}p{2.6cm}p{0.6cm}p{4.2cm}@{}}\toprule
Input & Source (comment) & Posted value & St. & Our value / note\\\midrule
GRH (conditional) & twinheterodox 698514; Stopple 698534 & $q \le 200$ & R & $d \le 199$; limit $64\pi^2e^{-2\gamma}$ \\
$7/8$ half-plane \cite{OAIqrh} & Tao 698471; LemmaLab 698472 & $655$; $653$ & R & $d_M = 653.98763\ldots$, $d \le 652$, kernel-checked \cite{CN1paper}\\
$11/12$ half-plane \cite{OAIqrh1112} & Tao 698471; LemmaLab 698473 & $1475$; $1151$ & R & $d \le 1151$ (edge form)\\
Banks \cite{Banks2026} & twinheterodox 698456; michaelmross 698466 & $10^{10}$ & R & $d < 10^{10}$, margin $6.7$\\
Banks + Lu--Zaman--Zhao & Larsson 698469 \cite{Larsson2026} & $50{,}000$ & C & \\
Inert-mass form & Silberman 698449/698451 & $10^{45}$ & C & source theorem unstated\\
Kerr \cite{Kerr2026} & Bryce 698475, who calls arXiv:2610.21099 (B.\ Kerr) ``my polynomial-lattice argument'' & $200{,}000$; $11{,}000$ & C & \\
Kimura \cite{Kimura2026} Prop.\ 9 & Rudolph 698480 & $1.89\cdot10^8$ & C & \\
Kimura chain thm 59 (Rudolph's input) & (ours) & --- & R & $d < 1.23\cdot10^8$ by \cref{prop:zero}\\
Kimura chain thm 58 & (ours) & --- & R & $d < 6.47\cdot 10^7$\\
Kimura Thm 1, parts 1/2 & (ours) & --- & R & $5.5\cdot10^9$; $1.61\cdot10^{15}$\\
Determinant, light & twinheterodox v1--v6.1 \cite{twinheterodox} & $\dots, 2000, 570$ & R & $570$ reproduced per row (\cref{sec:reproduction})\\
Determinant, elementary & twinheterodox 698521 & $2^{15}$ & R$^*$ & kernel-checked \cite{DETpaper} (comment 698632); external replay by Olivier (698668)\\
van Dijk \cite{vanDijk2026} & van Dijk 698529 & $8\cdot 10^9$ & C & \\
Green/van Dijk determinant via Kimi K3 & Pratt 698526 & Kimi $54{,}000$; a ChatGPT audit ``suggests'' a repair to $160{,}000$ & C & \\
van Dijk via Kimi K3 & Pratt 698682 & $<10^6$, ``basically fine'' & C & \\
Green, Thm 1.1 as stated \cite{Green2026} & --- & --- & --- & not yet extracted\\
OpenAI \cite{OAIsiegel} & --- & --- & --- & Lean theorem is existential\\
Goldfeld--Gross--Zagier--Oesterl\'e & Tao's post & $e^{7000}$; $e^{55}$ & C & classical\\
(unstated) & Anonymous 698759 & $187$ & NR & derivation not posted\\\bottomrule
\end{tabular}
\caption{Extractions of $D_0$. Values are quoted at source. ``R'' rows were recomputed by us in ball arithmetic and by a second, separately written implementation; R$^*$: a kernel-checked formalisation, not a recomputation; ``C'' rows are cited as posted.}\label{tab:survey}
\end{table}

\subsection{The Kimura rows}
Kimura's Lean development exports, besides Theorem~1, chain theorems \lean{chain_i : B4.GapFrom c_i l_i} meaning $(1-\beta)\log q \ge c_i$ for $\log q \ge l_i$ \cite{Kimura2026}. Through \cref{prop:zero}, \lean{chain_58} $= (2147/500000,\ 88711/5000)$ gives $d < 64{,}668{,}273$ and is the best of the $313$; \lean{chain_59}, the input of Rudolph's extraction, gives $d < 123{,}209{,}905$, against $1.89\cdot10^8$ by his Proposition~9 route. Both values are certified by us and rest on Kimura's Lean-formalised theorems, which we have not replayed. Rudolph's own value $1.89\cdot10^8$ comes from a different route (his Proposition~9 numerics), which we have not re-derived.

\subsection{One scale}\label{sec:scale}
\Cref{prop:edge} is a curve $\sigma_0 \mapsto d_M(\sigma_0)$, increasing on the sampled range from $199.12$ at $\sigma_0 \to \frac12$ to $653.99$ at $7/8$ and $1151.66$ at $11/12$. Reading it backwards assigns to any $D_0$ the abscissa $\sigma_{\mathrm{eq}}(D_0)$ that a real-zero-free segment would have to reach to give the same $D_0$. This is a calibration device, not a statement that a method proves anything about zeros.

\begin{table}[ht]\centering\small
\begin{tabular}{@{}lrr@{}}\toprule
Method & $D_0$ & $\sigma_{\mathrm{eq}}$\\\midrule
GRH & $199.12$ & $0.5$\\
determinant, light ($570$) & $570$ & $0.8610$\\
$7/8$ half-plane & $653.99$ & $0.8750$\\
$11/12$ half-plane & $1151.66$ & $0.9167$\\
determinant, elementary & $32{,}768$ & $0.9884$\\
Kimura chain 58 & $6.47\cdot10^7$ & $1 - 2.4\cdot10^{-4}$\\
Banks & $10^{10}$ & $1 - 1.9\cdot10^{-5}$\\\bottomrule
\end{tabular}
\caption{The equivalent zero-free abscissa. For large $D_0$, $1 - \sigma_{\mathrm{eq}} \approx (6/\pi)/\sqrt{D_0}$.}\label{tab:scale}
\end{table}

Two observations follow. On this problem the light determinant method is worth more than the $7/8$ half-plane. And within the zero-free family, a bound below $570$ would need a real-zero-free segment beginning at $0.8610$ uniformly for $q \ge 570$, which no unconditional result supplies. Classification-based routes (Weber's, as reconstructed in twinheterodox's historical companion, and Proposition~1.10 of \cite{BMOR2018}, whose proof for odd characters appeals to Watkins's tables) are outside the comparison.

%=====================================================================================================
\section{Reproducing the determinant bound \texorpdfstring{$q < 570$}{q < 570}}\label{sec:reproduction}
%=====================================================================================================
The bound is version 0.6.1 of twinheterodox \cite{twinheterodox}. In their own account it builds on Green's derivative determinant and Frobenius divisibility \cite{Green2026}, van Dijk's auxiliary-inert-prime nonvanishing \cite{vanDijk2026}, Pratt's suggestions (comments 698476, 698460), Kerr's higher Frobenius relations \cite{Kerr2026}, the Poisson expansion of michaelmross and Tao, and Rosser--Schoenfeld's prime estimates \cite{RS1962}. Their repository carries no licence; we use the mathematics, with citation, and our own words, code and tables.

\subsection{The argument in outline}
Suppose $h(K) = 1$ with $q$ in a band $[Q_0, Q_1]$. The argument compares two bounds for the same quantity, a smoothed sum over the noninert primes,
\[ T(X) = \sum_{p\ \text{noninert}} \frac{\log p}{p}\, f(p/X), \qquad f(u) = 2u\sum_{m\ge1}(1-mu)_+ , \]
which vanishes for $X \le q/4$.

\emph{The lower bound.} For a scale $U$, counting lattice points of $\OK$ in an annulus, with a Bonferroni sieve over a few small ideals, shows that the split primes cannot carry all the prime weight in $(U, \kappa U]$, so an inert prime $p_0$ lies there. With $R = p_0$, Green's derivative interpolation matrix, with rows $z^r w^s$ ($r < R$, $s < tR$) and columns the $k$-th $z$-derivatives ($k < R/9$) at the lattice points of a disk, has full row rank, by reduction modulo $p_0$, where Frobenius sends $z$ to $\bar z$ (van Dijk's argument). Fix a nonzero maximal minor $D$ and put $\Delta = |D|^2 \in \Z_{>0}$. At each inert prime $p$, triangular row changes by $z^p - w$, $z^{p^2} - z$ and $w^{p^2} - w$ make every entry of a row divisible by a power of $p$; summed over primes this is a lower bound for $\log\Delta$, whose main term is a sum over all primes minus $T(R)$. Factorials at the split primes and in the columns add further divisibility, again expressed through $T$ at smaller scales. An entrywise estimate bounds $\log\Delta$ from above. Comparing the two forces $T(R) \ge \alpha\log R - \frac\beta4\log q - K$, and, since $T$ is monotone, a piecewise affine lower bound $T(e^x) \ge L(x) - \gamma\log(q/Q_0)$ for $x$ across many scales, refined in eight stages that each feed the previous stage's bound back into the factorial gains.

\emph{The upper bound.} Under $h(K) = 1$, $\zeta_K(s) = \zeta(2s) + A_q(s)$ with $A_q$ the norm series of the principal form, so $-\frac{d}{ds}\log(\zeta_K(s)/\zeta(2s))$ is a Dirichlet series over the split primes and their powers. For a positive kernel $h(x) = Z^{-1}e^{-dx}P(e^{-ex})^2$, the average $\int h(x)\,T(e^x)\,dx$ is a combination of that series at the real points $s = 1 + d + je$, $j = 0, 1, \ldots$, which Poisson summation in one coordinate of the norm form evaluates up to an explicit error.

\emph{The contradiction.} If $\int h\,L > $ the upper bound at $q = Q_0$, and the upper bound decreases in $q$ faster than $\gamma\log(q/Q_0)$ across the band, then no $q$ in the band has $h(K) = 1$.

\subsection{The reproduction}
We derived each step independently and wrote certificate code that reads only v0.6.1's parameters. Our numbers were frozen and their checksums sent to the review session before any comparison. Three slips in our first code (two conservative, one not) were found by this comparison and fixed by derivation; none was an error in v0.6.1. With v0.6.1's parameters and its Mertens constant $1.334$ (the faithful configuration), our per-row envelope coefficients agree with theirs to within $6.3\cdot10^{-14}$ in each of the three bands we use, consistent with their outward rounding, and the lower bound of the band $[570, 580]$ is $10.9754131603100$ against their $10.975413160309$. The review session's separately written implementation agrees per row.

\subsection{Controls}
\begin{itemize}[leftmargin=*]
\item \textbf{Truth test.} At $q = 163$, $67$ and $43$, which have class number one, every inequality is a theorem. Our envelope lies below each field's actual noninert-prime sum, computed from its split primes up to $10^8$, at all nine stages (margin at least $1.9$). The Chowla--Selberg main term agrees with the actual weighted prime sum at $163$ to about $0.1\%$ (approximate: the tail beyond $10^8$ is estimated).
\item \textbf{Expected failure.} With v0.6.1's parameters the band $[560, 570]$ fails (gap $-0.026$), consistent with their record that kernel searches at $560$ ``did not yield a positive gap''.
\item \textbf{Degradation.} On the band $[565, 570]$, worsening an input at first order (Rosser--Schoenfeld's $E$ or Table~I by $1\%$, or one row's $\kappa$ by $1\%$) lowers the gap.
\item \textbf{Quadrature.} The kernel average computed through the cancelling exponential expansion agrees with a direct quadrature of the positive kernel to $2\cdot10^{-15}$.
\end{itemize}

%=====================================================================================================
\section{The improvements}\label{sec:improvement}
%=====================================================================================================
\begin{enumerate}[leftmargin=*]
\item \textbf{The Poisson modes.} In the norm series $A_q(s) = \sum_{b\ge1}\sum_{a\in\Z} ((a+b/2)^2 + qb^2/4)^{-s}$ the nonzero Poisson modes are
\[ A_q(s) - F_q(s) = \sum_{b,n\ge1} 4(-1)^{nb}\pi^s n^{s-1/2}(bv)^{1/2-s}K_{s-1/2}(2\pi n b v)/\Gamma(s),\qquad v = \sqrt q/2,\]
(the sign checked against lattice sums). Bounding them termwise with $K_\nu(x) \le \sqrt{\pi/(2x)}\,e^{-x+\nu^2/(2x)}$, sharp to $0.2\%$ here, instead of by a contour shift, reduces the error that multiplies the kernel's (large, cancelling) coefficients by a factor of $10^3$--$10^4$. This admits kernels of degree $20$ instead of $16$. It is load-bearing: with contour-shift accounting (our derivation of it, with $a_0 = 1$) the new band fails by $41.6$.
\item \textbf{The constant in $\sum_{p\le x}\log p/p$.} Rosser--Schoenfeld's (3.21) gives $\log x + E - 1/(2\log x)$ for all $x > 1$, with $E = -1.3325822757\ldots$; v0.6.1 uses $1.334$ for $x \ge 319$. This adds about $0.0016$ to the margin.
\item \textbf{Our own band data.} For each band we choose the auxiliary-prime rows (sieve sizes and the minimal $\kappa$ certified by the sieve inequality), the parameter $\gamma$ (only $\gamma \ge \beta/4$ is needed), and the kernel, as the top generalised eigenvector of the quadratic form that the kernel average defines, rounded to integers. All band data (rows, $\kappa$, sieve sizes, kernels, $\gamma$) are ours; the method constants $t = 3/64$, $C = 27/64$, $\tau = 3/10$ and the eight stages are v0.6.1's.
\item \textbf{Band uniformity.} The upper bound's derivative in $Y = \log(q/Q_0)$ is enclosed by a Taylor model whose coefficients are exact point evaluations, with a Cauchy-estimate remainder; the condition $U'(Y) < -\gamma$ makes the gap at $Q_0$ persist across the band.
\item \textbf{A certified constant.} $K(1/9) = \int_0^{1/9} I(k)\,dk$, for which v0.6.1 gives a closed form, is bracketed rigorously, $K(1/9) \in [0.0077989, 0.0078016]$, from certified values of $I(k)$ and its monotonicity (envelope theorem), and the upper end is used.
\end{enumerate}

%=====================================================================================================
\section{The theorem}\label{sec:theorem}
%=====================================================================================================
\begin{theorem}\label{thm:main}
If $h(\Q(\sqrt D)) = 1$ for a fundamental discriminant $D<0$, then $|D| < 565$. The argument excludes $[565, 652]$ by four certified bands of the determinant method and $|D| \ge 653$ by \cite{CN1paper}.
\end{theorem}

\begin{table}[ht]\centering\small
\begin{tabular}{@{}lcccc@{}}\toprule
Band & rows & kernel & gap at $Q_0$ & $U'$ on the band\\\midrule
$[565, 570]$ & 35 & deg.\ 20, $d=\frac14$, $e=\frac1{25}$ & $+0.002275$ & $\le -2.255$\\
$[570, 580]$ & 35 & deg.\ 20, $d=\frac14$, $e=\frac1{25}$ & $+0.020089$ & $\le -2.34$\\
$[580, 600]$ & 35 & deg.\ 20, $d=\frac14$, $e=\frac1{25}$ & $+0.055810$ & $\le -2.38$\\
$[600, 652]$ & 35 & deg.\ 20, $d=\frac14$, $e=\frac1{25}$ & $+0.124424$ & $\le -2.42$\\\bottomrule
\end{tabular}
\caption{The bands (our own band data; $\gamma = 2341/10000$ in each). The column $U'$ gives an upper bound over eight Taylor cells covering the band. Each band passes: a positive gap at $Q_0$; $U' < -\gamma$ on the band; and monotonicity of each Laplace component, so that the error bound at $Q_0$ is uniform.}\label{tab:bands}
\end{table}

\subsection{The finite range}\label{sec:finite}
For every fundamental discriminant with $3 \le q < 565$ (there are $173$) we compute $h$ directly as the number of reduced forms. It is $1$ exactly for $q \in \{3,4,7,8,11,19,43,67,163\}$.

\subsection{Size and fragility}
The cutoff moves from $570$ to $565$, about $0.9\%$. No candidate changes: the primes $q \equiv 19 \pmod{24}$ near the cut are $547$ and $571$, and both are excluded either way. The margin of the binding band $[565, 570]$, $0.0023$, is smaller than the effect of a $1\%$ change in Rosser--Schoenfeld's constant $E$ (about $0.015$), and the band certifies only with the Poisson-mode treatment of \cref{sec:improvement}; our attempts on $[560, 565]$, $[562, 565]$ and $[563, 565]$ did not certify.

%=====================================================================================================
\section{Limits}\label{sec:limits}
%=====================================================================================================
\begin{itemize}[leftmargin=*]
\item \textbf{Not kernel-checked.} The argument needs explicit prime estimates of Rosser--Schoenfeld quality at scales up to about $10^{16}$. The explicit-bounds development in PrimeNumberTheoremAnd (at commit \texttt{c39a7511}, 30 September 2026) formalises several such results but rests them on trusted numerical inputs; its Rosser--Schoenfeld files prove the asymptotic Mertens theorems, while their explicit statements are stated with \lean{sorry} (e.g.\ Theorem~12) or rest on files that contain \lean{sorry}. No three-axiom-clean source of these inputs exists in Lean yet. A kernel-checked version is future work.
\item \textbf{The survey is incomplete.} Rows marked C are cited, not checked.
\item \textbf{Two implementations, both AI-written.} The reproduction is by two separately written AI implementations (ball arithmetic, and high-precision floating point); no human has checked either.
\item \textbf{The inputs.} Rosser--Schoenfeld 1962 (Table I; Theorems 4, 6, 9, 14) is cited, read from the page images, and transcribed by us; OpenAI's $7/8$ theorem enters only through \cite{CN1paper} for $q \ge 653$.
\end{itemize}

%=====================================================================================================
\section{Credits, provenance and AI use}\label{sec:credits}
%=====================================================================================================
The determinant method and the bound $570$ are twinheterodox's \cite{twinheterodox}, with the credits given there. The zero-forcing route is twinheterodox's and michaelmross's; the evaluation at the edge is LemmaLab's; the $11/12$ value $1151$ is LemmaLab's; the Kimura extraction route is Rudolph's; the inert-prime facts are classical, as emphasised by Silberman and Pratt. Our contributions are the survey and its verification, the strength scale, the reproduction, and the improvements listed in \cref{sec:improvement}.

This work was carried out with Claude (Anthropic) in Claude Code: derivations, code and text were AI-generated, and every result was reviewed by a separate AI session that wrote its own implementation. These are AI reviews, not human peer review. The author directed the work and takes responsibility for the claims.

\bibliographystyle{plain}
\bibliography{refs}
\end{document}
